\chapter*{Appendix K: Homotopy and Derived Methods}
\addcontentsline{toc}{chapter}{Appendix K: Homotopy and Derived Methods}

\section{Basic Homotopy Structures}

\subsection{Paths and Homotopies}

Let $X$ be a topological space.  
A path in $X$ is a continuous map $\gamma : [0,1] \to X$.  
Two paths $\gamma_0,\gamma_1$ with the same endpoints are homotopic if there exists a map
\[
H : [0,1]\times[0,1] \to X,
\qquad
H(0,t)=\gamma_0(t),\;
H(1,t)=\gamma_1(t),
\]
with fixed endpoints.

\begin{definition}[Fundamental Group]
The fundamental group $\pi_1(X,x_0)$ is the set of homotopy classes of loops based at $x_0$, with concatenation as the group operation.
\end{definition}

\subsection{Higher Homotopy Groups}

For $n\ge 1$,
\[
\pi_n(X,x_0) = [S^n,X],
\]
the set of homotopy classes of based maps $S^n\to X$.  
For $n\ge 2$, $\pi_n(X)$ is abelian.

\section{Simplicial Objects}

\subsection{Simplicial Sets}

A simplicial set is a functor
\[
X_\bullet : \Delta^{\mathrm{op}} \to \mathbf{Set},
\]
with face and degeneracy maps satisfying the simplicial identities.

\subsection{Geometric Realisation}

The geometric realisation of a simplicial set $X_\bullet$ is
\[
|X_\bullet|
  = \left( \bigsqcup_{n\ge 0} X_n \times \Delta^n \right) / \sim.
\]

\section{Chain Complexes and Derived Categories}

\subsection{Chain Complexes}

A chain complex $(C_\bullet,d_\bullet)$ is a family of abelian groups with boundary maps
\[
\cdots \xrightarrow{d_3} C_2 \xrightarrow{d_2} C_1 \xrightarrow{d_1} C_0 \to 0
\]
satisfying $d_{n} \circ d_{n+1} = 0$.

\subsection{Chain Maps and Homotopies}

Chain maps commute with boundaries:
\[
f_{n-1}\circ d_n = d'_n\circ f_n.
\]
Chain-homotopic maps satisfy
\[
f - g = d'H + Hd.
\]

\subsection{Derived Category}

The derived category $D(A)$ of an abelian category $A$ is obtained by formally inverting quasi-isomorphisms of chain complexes:
\[
\text{quasi-isomorphisms} :
H_\bullet(f) \text{ isomorphism}.
\]

\section{Model Categories}

\subsection{Definitions}

A model category $(\mathcal{M},\mathsf{W},\mathsf{C},\mathsf{F})$ consists of:
\begin{itemize}
\item weak equivalences $\mathsf{W}$,
\item cofibrations $\mathsf{C}$,
\item fibrations $\mathsf{F}$,
\end{itemize}
satisfying the Quillen axioms:
2-out-of-3, retract stability, lifting conditions, and factorisation conditions.

\subsection{Homotopy Limit and Colimit}

Homotopy limits and colimits are defined using cofibrant/fibrant replacements:
\[
\operatorname{holim} X \simeq \lim R(X),
\qquad
\operatorname{hocolim} X \simeq \operatorname{colim} Q(X).
\]

\section{\texorpdfstring{$\infty$}{∞}-Categories}

\subsection{Quasi-Categories}

A quasi-category is a simplicial set $C_\bullet$ in which every inner horn $\Lambda^n_k \to C$ admits a filler $\Delta^n\to C$.

\subsection{Mapping Spaces}

For objects $x,y$ in a quasi-category $C$, the mapping space is
\[
\mathrm{Map}_C(x,y)
  = \{ \text{1-simplices with $x$ as source and $y$ as target} \}
\]
with higher simplices given by homotopies.

\section{Homotopy-Coherent Diagrams}

\subsection{Coherent Natural Transformations}

Given functors $F,G: I \to \mathcal{C}$ into an $\infty$-category, a homotopy-coherent natural transformation consists of:
\[
\{\eta_i : F(i)\to G(i)\}
\]
together with higher homotopies relating $\eta$ to the compositions along each morphism in $I$.

\subsection{Coherent Limits}

The homotopy limit of a diagram $D:I\to \mathcal{C}$ is the universal cone in which all commutativity conditions hold up to coherent homotopy.

\section{Spectral Sequences}

\subsection{Filtered Complexes}

Let $(C_\bullet,F_\bullet)$ be a filtered complex with
\[
0 \subseteq F_0 C_n \subseteq F_1 C_n \subseteq \cdots \subseteq C_n.
\]

\subsection{$E_r$-Pages}

The associated spectral sequence has pages
\[
E^1_{p,q} = H_{p+q}(F_p C / F_{p-1}C),
\qquad
d^r : E^r_{p,q} \to E^r_{p-r,q+r-1}.
\]

\section{Homotopy Fiber and Cofiber}

\subsection{Fiber}

For a map $f:X\to Y$, the homotopy fiber $\operatorname{hofib}(f)$ is
\[
\operatorname{hofib}(f)
  = \{(x,\gamma) : x\in X,\; \gamma:[0,1]\to Y,\; \gamma(0)=f(x),\gamma(1)=y_0\}.
\]

\subsection{Cofiber}

The homotopy cofiber is defined as
\[
\operatorname{hocofib}(f)
  = Y \cup_f (X\times[0,1]) / (X\times\{1\}),
\]
modding out the top cylinder.

\section{Homotopy Kan Extensions}

\subsection{Left Extension}

Given $F: I\to\mathcal{C}$ and $u:I\to J$, the left homotopy Kan extension is
\[
(\operatorname{Lan}_u F)(j)
  = \operatorname{hocolim}_{(i\downarrow j)} F(i).
\]

\subsection{Right Extension}

Similarly,
\[
(\operatorname{Ran}_u F)(j)
  = \operatorname{holim}_{(j\downarrow i)} F(i).
\]

\section{Derived Functors}

\subsection{Left and Right Derived Functors}

For functor $F: A\to B$,
\[
\mathbf{L}F(X) = F(QX),
\qquad
\mathbf{R}F(X) = F(RX),
\]
using cofibrant or fibrant replacements.

\subsection{Tor and Ext}

In module categories:
\[
\operatorname{Tor}^A_n(M,N) = H_n(P_\bullet \otimes_A N),
\qquad
\operatorname{Ext}^n_A(M,N) = H^n(\operatorname{Hom}(P_\bullet,N)),
\]
where $P_\bullet$ is a projective resolution.


